
\section{Preliminaries}
\label{section:preliminaries}\label{sec:preliminaries}

\paragraph{General notation.}{ We use $\A_{i,j}$ to denote the $(i,j)$-th entry of $\A$. For $\A \in \R^{n\times n}$; we use $\Spectrum{\A}$ for its spectrum; $\lambda_i(\A)$ and $\sigma_i(\A)$ for its $i$-th smallest eigenvalue and singular value respectively; and $\rho(\A) := \abs{\lambda_n(\A)}$ for its spectral radius. We use $\Vert \cdot\Vert_p$ for the $\ell_p$-norm. When $\A$ is PSD, $\lambdamin(\A)$ denotes its smallest nonzero eigenvalue and $\kappa(\A):=\lambda_n(\A)/\lambdamin(\A)$ denotes its pseudo-condition number. We use $\langle \cdot, \cdot \rangle$ for the Euclidean inner product; $\one$ for the all ones vector; and $e_i$ for the $i$-th standard basis vector. We define $\deltaij := e_i - e_j$ and $[k] \defeq \{1, ..., k\}$. We use $x \approx_\epsilon y$ as shorthand for $(1-\epsilon) y \leq x \leq (1+\epsilon) y$. For $v \in \R^n$, we use $v[i:j]$ for the sub-vector from index $i$ to $j$.}

\paragraph{Graphs.}{We use $G = (V, E, w)$ to denote a weighted undirected graph on $V$ with edges $E$ and edge weights $w \in \R^{E}_{> 0}$ (or $G = (V, E)$ if unweighted). We use $\A_G$ to denote its (weighted) adjacency matrix $(\A_G)_{u,v} = w_{u,v}$ for $u,v \in V \times V$ and $\D_G$ to denote its diagonal (weighted) degree matrix $(\D_G)_{u} = \sum_{\{u, v\} \in E} w_{u,v}$ for $u \in V$ (treated as $w_{u,v} = 1$ if $G$ is unweighted). We define $\laplacian{G} := \D_G - \A_G$ as its graph Laplacian. $\dmax(G)$ and ${\dmin}(G)$ refer to the max and min diagonal entry in $\D_G$. We may drop the argument or subscript $G$ it is clear from context. The effective resistance between nodes $i, j \in V$ is denoted $\rGij = \deltaij^\top \laplacian{G}^\dagger \deltaij$. We assume all input graphs are connected, as effective resistances can be computed separately on connected components.
}

\paragraph{Symmetric Diagonally Dominant (SDD) Matrices}{
A matrix $\M \in \R^{n \times n}$ is SDD if it can be decomposed as $\M = \diagonal{\M} - \adjacency{\M}$, where the $\diagonal{\M}$ is a diagonal matrix with non-negative entries and $\adjacency{M}$ is a matrix with zeros on the diagonal such that $d_{i,i} > \sum_{j = 1}^n |a_{i,j}|$. We define the normalization of $\M$ as $\normalized{M} := \diagonal{M}^{-1/2}\M \diagonal{M}^{-1/2}$. Throughout this paper, we assume, without loss of generality that $\D_M$ has strictly positive entries on the diagonal (otherwise, we can simply remove an entire row and column of zeros). We use ${\dmax}(\M)$ and ${\dmin}(\M)$ to denote the max and min entry in the diagonal of $\D_M$ respectively. We may drop the argument or subscript $\M$ if it is clear from context. }

\paragraph{Runtimes and Space Complexities.} {In our algorithmic results and analysis, when clear from context, we use $\widetilde{O}(\cdot)$ notation to hide polylogarithmic factors in the number of vertices, the number of edges, the size of the matrix, the number of nonzero entries in a matrix, the maximum and minimum diagonal element of a matrix, the maximum and minimum weighted degree of a graph, $\epsilon$, the condition number, and the normalized psuedo-condition number of a matrix (see Definition~\ref{def:pseudo-condition-number}). We say event $E$ occurs with high probability in $t$ if $\Prob{E} \geq 1 - t^{-c}$, where $c> 0$ can be controlled by appropriately configuring the algorithm parameters. We may simply say that an event occurs ``with high probability'' or ``whp.'' if it occurs with high probability in the dimension of a matrix or number of nodes in a graph.} 